\input{template.tex}

\ifPDFTeX
  \renewcommand{\rmdefault}{cmr}
\fi

\def\SheetNo{5}
\setlength{\emergencystretch}{2em}

\begin{document}
\sheettitle{Комбинаторика.}
\begin{center}\textbf{Решения}\end{center}

\begin{problems}

\item Сколькими способами можно распределить $15$ различных гостей по четырём различным комнатам, заняв каждую комнату хотя бы одним гостем?
\par\smallskip\noindent\textbf{Решение.} Без ограничений имеется $4^{15}$ расселений: каждому гостю независимо выбирается одна из четырёх комнат. Для заранее выбранных $j$ пустых комнат число расселений в остальные равно $(4-j)^{15}$, а выбрать эти комнаты можно $C_4^j$ способами. По формуле включений-исключений ответ равен
\[
4^{15}-4\cdot3^{15}+6\cdot2^{15}-4\cdot1^{15}.
\]\par

\item Имеется $r$ различных предметов и $n+p$ человек. Докажите, что число распределений предметов, при которых каждый из заранее указанных $n$ человек получает хотя бы один предмет, равно
\[
\sum_{j=0}^{n}(-1)^j C_n^j(n+p-j)^r.
\]
\par\smallskip\noindent\textbf{Решение.} Без ограничений каждый из $r$ предметов можно отдать любому из $n+p$ человек, что даёт $(n+p)^r$ распределений. Для каждого из указанных $n$ человек рассмотрим событие, что он не получил ни одного предмета. Если запрещено отдавать предметы заранее выбранным $j$ из этих людей, остаётся $(n+p-j)^r$ распределений. Выбрать этих людей можно $C_n^j$ способами. Применяя включения-исключения, получаем
\[
\sum_{j=0}^{n}(-1)^j C_n^j(n+p-j)^r.
\]\par

\item Сколькими способами можно заполнить таблицу $k\times l$ числами $+1$ и $-1$, чтобы произведение в каждой строке и в каждом столбце было равно $1$?
\par\smallskip\noindent\textbf{Решение.} Ответ: $2^{(k-1)(l-1)}$. Произвольно заполним первые $k-1$ строк и $l-1$ столбцов. Для каждой из этих строк последняя клетка определяется однозначно условием, что произведение равно $1$; аналогично определяются клетки последней строки в первых $l-1$ столбцах.
\par Значение правой нижней клетки, необходимое для последней строки, равно произведению всех элементов исходного прямоугольника $(k-1)\times(l-1)$. Условие для последнего столбца даёт то же произведение. Поэтому противоречия нет и каждое начальное заполнение имеет ровно одно продолжение.\par

\item Каково наибольшее возможное число различных точек пересечения диагоналей $n$-угольника?
\par\smallskip\noindent\textbf{Решение.} При $n=3$ пересечений нет. Пусть $n\ge4$. Рассмотрим выпуклый многоугольник, никакие три диагонали которого не проходят через одну внутреннюю точку. Такого положения можно добиться малым изменением положений вершин.
\par Зафиксируем диагональ. Пусть на меньшей из двух дуг между её концами находится $k$ других вершин, а на большей - $n-k-2$. Тогда
\[
1\le k\le\left\lfloor\frac{n-2}{2}\right\rfloor.
\]
Диагональ пересекается с отрезками, соединяющими вершину с одной стороны с вершиной с другой, то есть имеет $k(n-k-2)$ различных внутренних точек пересечения.
\par Пронумеруем вершины против часовой стрелки числами $0,\ldots,n-1$. Для фиксированного $k$ проведём диагонали вида $\{x,x+k+1\}$, считая номера по модулю $n$. Две записи могут задать одну диагональ лишь при перестановке концов; для этого необходимо и достаточно
\[
x+2k+2\equiv x\pmod n.
\]
В указанном диапазоне $k$ это возможно только при чётном $n$ и $k=n/2-1$. Поэтому обычно диагоналей данного типа $n$, а в исключительном случае - $n/2$, поскольку каждая считается с двух концов.
\par Просуммируем количества пересечений по всем диагоналям и разделим на $2$: каждая точка учитывается на обеих пересекающихся диагоналях.
\par\textit{Нечётное $n=2m+1$.} Возможны $k=1,\ldots,m-1$, каждого типа имеется $n$ диагоналей. Число пересечений равно
\[
\frac{2m+1}{2}\sum_{k=1}^{m-1}k(2m-1-k)
=\frac{(2m+1)m(m-1)(2m-1)}{6}
=C_n^4.
\]
\textit{Чётное $n=2m$.} Для $k=1,\ldots,m-2$ имеется $n$ диагоналей, а для $k=m-1$ - $m$. Поэтому число пересечений равно
\[
\begin{aligned}
\frac12\left(2m\sum_{k=1}^{m-2}k(2m-2-k)+m(m-1)^2\right)
&=\frac{m(m-1)(2m-1)(2m-3)}{6}\\
&=C_n^4.
\end{aligned}
\]
Обе суммы вычисляются по формулам для суммы первых натуральных чисел и суммы их квадратов.
\par Наконец, при любом расположении вершин каждая внутренняя точка пересечения задаётся парой диагоналей с четырьмя различными концами. Для одной четвёрки вершин внутри отрезков может пересекаться не более одной пары соединяющих их отрезков. Следовательно, различных точек не больше $C_n^4$. Построенный выпуклый многоугольник достигает этой оценки. Ответ:
\[
\boxed{C_n^4=\frac{n(n-1)(n-2)(n-3)}{24}}.
\]\par
\par\textit{Короткий способ.} Возьмём выпуклый $n$-угольник, в котором никакие три диагонали не пересекаются в одной внутренней точке. Любые четыре вершины образуют выпуклый четырёхугольник, диагонали которого дают ровно одну внутреннюю точку пересечения. Обратно, каждая точка пересечения двух диагоналей определяет четвёрку их концов. Получаем биекцию между четвёрками вершин и точками пересечения, поэтому последних $C_n^4$.
\par Больше точек получить нельзя при любом расположении вершин: для каждой четвёрки существует не более одной пары соединяющих её вершины отрезков, пересекающихся во внутренних точках. Следовательно, максимум равен $C_n^4$.\par

\item На выборах один кандидат получил $12$ голосов, а другой - $8$. Сколько существует порядков подсчёта голосов, при которых после каждого подсчитанного голоса первый кандидат строго опережает второго? Порядки различаются последовательностью голосов за двух кандидатов.
\par\smallskip\noindent\textbf{Решение.} Голосу за первого кандидата сопоставим шаг $(1,1)$, а за второго - $(1,-1)$. Получим путь из $(0,0)$ в $(20,4)$, который после первого шага должен находиться строго выше оси абсцисс.
\par Первый шаг обязательно направлен вверх. Удалим его и перенесём начало оставшегося пути в $(0,0)$. Требуется посчитать пути из $(0,0)$ в $(19,3)$, не достигающие прямой $y=-1$. Всего путей в $(19,3)$ имеется $C_{19}^{8}$: в каждом $11$ шагов вверх и $8$ вниз.
\par Назовём путь плохим, если он достигает $y=-1$. Пусть $(k,-1)$ - первая точка достижения этой прямой. Отразим участок пути после неё относительно $y=-1$. Начало и длина пути сохраняются, а конец переходит из $(19,3)$ в $(19,-5)$.
\par Это биекция между плохими путями и всеми путями из $(0,0)$ в $(19,-5)$: любой путь с таким концом обязательно достигает $y=-1$, и отражение его хвоста после первого достижения возвращает исходный плохой путь. В пути к $(19,-5)$ имеется $7$ шагов вверх и $12$ вниз, поэтому плохих путей $C_{19}^{7}$. Ответ равен
\[
C_{19}^{8}-C_{19}^{7}=75582-50388=25194.
\]\par

\item Пусть $E_n$ - число перестановок чисел $1,\ldots,n$, в которых ровно один раз следующее число больше предыдущего. Например, такова перестановка $4312$. Докажите, что при каждом натуральном $n$
\[
E_{n+1}=2E_n+n.
\]
\par\smallskip\noindent\textbf{Решение.} Рассмотрим вставку числа $n+1$ в перестановку чисел $1,\ldots,n$. Подъёмом будем называть пару соседних элементов $a_i<a_{i+1}$. При вставке наибольшего числа количество подъёмов не уменьшается, поэтому перестановки с двумя и более подъёмами не подходят.
\par Если подъёмов нет, перестановка единственна: $n,n-1,\ldots,1$. Вставка $n+1$ в любую позицию, кроме начала, создаёт ровно один подъём. Таких позиций $n$.
\par Если подъём ровно один, можно вставить $n+1$ в начало: нового подъёма не возникнет. Также можно вставить его между элементами единственного подъёма: пара $a_i<a_{i+1}$ заменится на $a_i<n+1>a_{i+1}$, и подъём останется один. Любая другая вставка добавляет второй подъём. Значит, каждая из $E_n$ перестановок даёт ровно два варианта.
\par Все полученные перестановки различны, поскольку исходная однозначно восстанавливается удалением числа $n+1$. Итого $E_{n+1}=n+2E_n$.\par

\item Сколько целых чисел в промежутке от $1$ до $999$ включительно не делятся ни на $3$, ни на $5$, ни на $7$?
\par\smallskip\noindent\textbf{Решение.} По формуле включений-исключений искомое число равно
\[
\begin{aligned}
999
&-\left\lfloor\frac{999}{3}\right\rfloor
-\left\lfloor\frac{999}{5}\right\rfloor
-\left\lfloor\frac{999}{7}\right\rfloor\\
&+\left\lfloor\frac{999}{15}\right\rfloor
+\left\lfloor\frac{999}{21}\right\rfloor
+\left\lfloor\frac{999}{35}\right\rfloor
-\left\lfloor\frac{999}{105}\right\rfloor
=457.
\end{aligned}
\]\par

\item Пусть $p$ - простое число и $0<k<p$. Докажите, что $p\mid C_p^k$:
\begin{parts}
  \item с помощью формулы для биномиального коэффициента;
  \item комбинаторно, без использования формул.
\end{parts}
\par\smallskip\noindent\textbf{Решение.} \textbf{а)} По формуле для биномиального коэффициента
\[
k!\,C_p^k=p(p-1)\cdots(p-k+1).
\]
Правая часть делится на $p$, а $\gcd(k!,p)=1$, поскольку $k<p$ и $p$ простое. Следовательно, $p\mid C_p^k$.
\par\textbf{б)} Расположим $p$ элементов на окружности в равноотстоящих позициях. Каждому $k$-элементному подмножеству сопоставим раскраску, в которой ровно его позиции красные. Две раскраски назовём эквивалентными, если одна получается из другой поворотом на целое число позиций. Это отношение разбивает все $k$-элементные подмножества на классы.
\par Докажем, что в каждом классе ровно $p$ подмножеств. Рассмотрим все повороты одной раскраски на $0,1,\ldots,p-1$ позиций. Если два результата совпали, исходная раскраска сохраняется при некотором сдвиге на $t$ позиций, где $0<t<p$. Поскольку $p$ простое, последовательные сдвиги на $t$ обходят все $p$ позиций. Значит, все позиции должны иметь одинаковый цвет. Это невозможно при $0<k<p$. Поэтому все $p$ поворотов различны.
\par Следовательно, множество всех $k$-элементных подмножеств разбито на классы по $p$ элементов, и его мощность $C_p^k$ делится на $p$.\par

\item На полке стоят $n$ книг. Сколькими способами можно выбрать $p$ из них, если между соседними выбранными книгами и после последней выбранной должно оставаться хотя бы $s$ книг? Считайте, что $n\ge(s+1)p$.
\par\smallskip\noindent\textbf{Решение.} Пусть $x_0$ - число книг до первой выбранной, $x_i$ при $1\le i<p$ - число книг между $i$-й и $(i+1)$-й выбранными, а $x_p$ - число книг после последней выбранной. Тогда
\[
x_0+x_1+\cdots+x_p=n-p,\qquad x_0\ge0,\quad x_i\ge s\quad(1\le i\le p).
\]
Положим $y_0=x_0$ и $y_i=x_i-s$ при $1\le i\le p$. Получим уравнение
\[
y_0+y_1+\cdots+y_p=n-p(s+1)
\]
в неотрицательных целых числах. Каждый его набор решений однозначно задаёт выбранные книги. По методу шаров и перегородок число решений равно
\[
C_{n-p(s+1)+p}^{p}=C_{n-ps}^{p}.
\]\par

\item Найдите коэффициент при $x^{10}y^{10}z^{10}$ в многочлене
\[
(x+y)^{10}(y+z)^{10}(z+x)^{10}.
\]
\par\smallskip\noindent\textbf{Решение.} Если из первого множителя выбрано слагаемое $x^ky^{10-k}$, то из второго необходимо выбрать $y^kz^{10-k}$, а из третьего - $z^kx^{10-k}$. Коэффициент при произведении этих слагаемых равен $(C_{10}^k)^3$. Поэтому искомый коэффициент равен
\[
\sum_{k=0}^{10}(C_{10}^k)^3.
\]\par

\item Докажите, что каждое из следующих количеств равно $n$-му числу Каталана $c_n$:
\begin{parts}
  \item число заполнений таблицы $2\times n$ числами $1,\ldots,2n$, каждое из которых используется один раз, с возрастанием слева направо в строках и сверху вниз в столбцах;
  \item число способов соединить $2n$ точек окружности попарно $n$ хордами так, чтобы хорды не пересекались;
  \item число параллеломино периметра $2n+2$. Параллеломино задаётся парой ломаных по линиям клетчатой сетки: они начинаются в $(0,0)$, имеют общий конец, идут только вверх и вправо и не имеют других общих точек. Периметр - сумма длин двух ломаных.
\end{parts}
\par\smallskip\noindent\textbf{Решение.} \textbf{а)} Просмотрим числа $1,2,\ldots,2n$ по возрастанию. Если число стоит в верхней строке, запишем открывающую скобку; если в нижней - закрывающую. Получим правильную скобочную последовательность: в любом начальном отрезке чисел нижних элементов не больше, чем верхних, поскольку над каждым нижним элементом стоит меньший верхний.
\par Обратно, по правильной скобочной последовательности поместим номера открывающих скобок в верхнюю строку, а закрывающих - в нижнюю, в обоих случаях по возрастанию. Условие на начальные отрезки гарантирует возрастание в столбцах. Эти преобразования взаимно обратны, поэтому таблиц столько же, сколько правильных скобочных последовательностей из $n$ пар скобок, то есть $c_n$.
\par\textbf{б)} Докажем утверждение по индукции. При $n=0$ имеется одно пустое разбиение, что соответствует $c_0=1$. Пронумеруем $2n$ точек против часовой стрелки числами $0,\ldots,2n-1$ и рассмотрим хорду из точки $0$. По обе стороны от неё должно остаться чётное число точек, поэтому второй конец имеет вид $2k-1$, где $1\le k\le n$.
\par Хорда $(0,2k-1)$ оставляет с одной стороны $2k-2$ точки, а с другой - $2n-2k$. Остальные хорды не могут соединять эти две части, иначе они пересекли бы выбранную хорду. По предположению индукции части можно независимо разбить $c_{k-1}$ и $c_{n-k}$ способами, что даёт $c_{k-1}c_{n-k}$ вариантов.
\par Случаи для разных $k$ не пересекаются: в одном разбиении из точки $0$ выходит ровно одна хорда, поэтому два разных её конца невозможны. Каждый допустимый вариант попадает в один из этих случаев. Значит, общее число равно
\[
\sum_{k=1}^{n}c_{k-1}c_{n-k}=c_n
\]
по рекуррентному соотношению для чисел Каталана.
\par\textbf{в)} Построим биекцию с правильными скобочными последовательностями из $n$ пар скобок. Обозначим шаг вверх через $U$, вправо - через $R$. Каждая из двух ломаных имеет длину $n+1$. Будем проходить их одновременно: после одинакового числа $k$ шагов обе точки находятся на прямой $x+y=k$.
\par Назовём ломаные верхней и нижней по их взаимному расположению. Верхняя начинается шагом $U$ и заканчивается шагом $R$, нижняя начинается шагом $R$ и заканчивается шагом $U$. Удалим эти начальные и конечные шаги; в каждой ломаной останется $n-1$ шаг.
\par Каждую пару соответствующих оставшихся шагов закодируем двумя скобками:
\begin{center}
\begin{tabular}{c|c|c}
Шаг верхней & Шаг нижней & Скобки\\\hline
$U$ & $R$ & \texttt{((}\\
$R$ & $U$ & \texttt{))}\\
$U$ & $U$ & \texttt{()}\\
$R$ & $R$ & \texttt{)(}
\end{tabular}
\end{center}
В начало записи добавим одну открывающую скобку, в конец - одну закрывающую. Всего получится $2(n-1)+2=2n$ скобок.
\par Пусть $\Delta$ - разность высот верхней и нижней ломаных после одинакового числа шагов. Между общими началом и концом всегда $\Delta\ge1$, иначе ломаные встретились бы. После первых шагов $\Delta=1$, а баланс начальной открывающей скобки равен $1$. Таблица поддерживает равенство
\[
\text{баланс на границе пар}=2\Delta-1:
\]
первые две строки меняют $\Delta$ на $+1$ или $-1$, а баланс - на $+2$ или $-2$; последние две строки ничего не меняют. Поэтому на границах пар баланс не меньше $1$, а внутри пары он может уменьшиться не более чем на единицу и не становится отрицательным. Перед последними шагами $\Delta=1$, так что заключительная закрывающая скобка доводит баланс до нуля. Получена правильная скобочная последовательность.
\par Обратно, у любой правильной скобочной последовательности длины $2n$ удалим первую и последнюю скобки, разобьём оставшуюся запись на пары и восстановим шаги по таблице. Добавим вынужденные начальные и конечные шаги. На границах пар длина префикса исходной записи нечётна; его баланс неотрицателен и нечётен, значит, положителен. Следовательно, $\Delta\ge1$, и ломаные не встречаются до последнего шага. Перед последними шагами баланс равен $1$, поэтому $\Delta=1$, а после них $\Delta=0$. При равном числе шагов это означает общий конец. Получена фигура периметра $2n+2$.
\par Построения взаимно обратны, поэтому число параллеломино равно $c_n$.\par

\item Докажите, что число корневых упорядоченных полных бинарных деревьев с $2n+1$ вершиной равно числу Каталана $c_n$. Здесь у каждой вершины либо нет потомков, либо есть ровно два - левый и правый; порядок потомков учитывается. Постройте явную биекцию с правильными скобочными последовательностями из $n$ пар скобок.
\par\smallskip\noindent\textbf{Решение.} Обойдём дерево поиском в глубину, всегда сначала посещая левое поддерево, затем правое. При переходе от вершины к левому потомку запишем открывающую скобку, а при переходе к правому - закрывающую. Возвращения к родителю ничего не добавляют.
\par Обозначим полученную строку через $S(T)$. Для дерева из одного листа она пуста. Если у корня имеются поддеревья $T_L,T_R$, то
\[
S(T)=\texttt{(}\,S(T_L)\,\texttt{)}\,S(T_R).
\]
По индукции эта строка является правильной скобочной последовательностью. В полном бинарном дереве с $2n+1$ вершиной ровно $n$ внутренних вершин: каждая даёт две исходящие дуги, а всего дуг $2n$. Каждой внутренней вершине соответствует одна пара скобок, поэтому в строке $n$ пар.
\par Обратное построение также задаётся рекурсивно. Пустой строке соответствует лист. Непустая правильная скобочная последовательность единственным образом представляется в виде $\texttt{(}\,u\,\texttt{)}\,v$, где указанная закрывающая скобка парная первой, а $u,v$ - правильные скобочные последовательности. Построим по $u$ левое поддерево, по $v$ - правое и присоединим их к новому корню. Построения взаимно обратны. Следовательно, искомое число деревьев равно $c_n$.\par

\item На плоскости выбраны $n$ точек: никакие три не лежат на одной прямой, а никакие четыре - на одной окружности. Проведены все прямые через пары выбранных точек и все окружности через тройки. Каково наибольшее возможное число различных точек пересечения этих прямых с этими окружностями?
\par\smallskip\noindent\textbf{Решение.} При $n<3$ окружностей нет, поэтому ответ $0$. Пусть $n\ge3$. Всего проведено $C_n^2$ различных прямых и $C_n^3$ различных окружностей. Зафиксируем прямую через исходные точки $A,B$ и разделим окружности на три группы.
\par Через обе точки $A,B$ проходят $n-2$ окружностей. Все они пересекают выбранную прямую ровно в этих двух точках, поэтому дают две различные точки пересечения.
\par Ровно через одну из точек $A,B$ проходят $2C_{n-2}^2$ окружностей: выбираем одну из этих точек и ещё две из остальных. Каждая такая окружность даёт не более одной новой точки пересечения.
\par Ни через $A$, ни через $B$ не проходят $C_{n-2}^3$ окружностей. Каждая даёт не более двух точек пересечения. Следовательно, на выбранной прямой имеется не более
\[
2+2C_{n-2}^2+2C_{n-2}^3
\]
различных точек пересечения с окружностями. Здесь биномиальные коэффициенты с нижним индексом меньше верхнего считаются равными нулю.
\par Просуммируем оценку по всем $C_n^2$ прямым. Каждая исходная точка учтена $n-1$ раз, поэтому для неё нужно убрать $n-2$ лишних учётов. Другие возможные совпадения только уменьшают число различных точек. Получаем верхнюю оценку
\[
\begin{aligned}
C_n^2\left(2+2C_{n-2}^2+2C_{n-2}^3\right)-n(n-2)
&=n+n(n-1)C_{n-1}^3\\
&=n+\frac{n(n-1)^2(n-2)(n-3)}{6}.
\end{aligned}
\]
\par Покажем достижимость. Сначала расположим все $n$ точек на одной окружности. Окружности через тройки тогда совпадают с ней, а каждая прямая через пару исходных точек пересекает её в двух различных точках и не является касательной.
\par Достаточно малое возмущение исходных точек сохраняет два различных пересечения каждой прямой с каждой окружностью. Это можно обеспечить одновременно для всех пар, поскольку их конечное число, а окружность через три различные исходные точки непрерывно зависит от их положения в окрестности выбранной конфигурации.
\par Выберем возмущение в общем положении: без трёх исходных точек на прямой, четырёх на окружности и без дополнительных совпадений новых точек пересечения. Такие совпадения задают специальные соотношения между координатами, которых можно избежать при сколь угодно малом общем возмущении. Тогда на каждой прямой достигается указанное число точек, а единственными повторными учётами между разными прямыми остаются исходные точки. Следовательно, оценка достигается. Ответ при $n\ge3$:
\[
\boxed{n+\frac{n(n-1)^2(n-2)(n-3)}{6}}.
\]\par

\item Во всех клетках доски $2020\times2020$ расставлены буквы В, С, О и Ш. Расстановка называется гармоничной, если в каждом квадрате $2\times2$ все четыре буквы различны. Найдите число гармоничных расстановок.
\par\smallskip\noindent\textbf{Решение.} Соседние клетки каждой строки содержат разные буквы, поскольку входят в некоторый квадрат $2\times2$. Если в каждой тройке последовательных клеток первая и третья буквы совпадают, строка чередует две буквы. Поэтому, если в строке встречаются хотя бы три разные буквы, найдётся тройка подряд с попарно различными буквами.
\par Пусть такая тройка имеет вид $A,B,C$, а четвёртая буква - $D$. В соседней строке под ней обязана стоять тройка $C,D,A$: средняя буква должна дополнять до четырёх букв одновременно пары $A,B$ и $B,C$, поэтому это $D$, после чего крайние буквы определяются однозначно. Следующая строка снова содержит $A,B,C$. Продолжая в обе стороны, получаем
\[
\begin{array}{ccc}
A&B&C\\
C&D&A\\
A&B&C\\
C&D&A\\
\vdots&\vdots&\vdots
\end{array}
\]
Таким образом, эти три столбца чередуют по две буквы. Если один столбец чередует две буквы, то в соседнем столбце каждая пара соседних клеток должна содержать две оставшиеся буквы. Значит, соседний столбец также чередует две буквы. Продолжая влево и вправо, заключаем, что все столбцы доски чередуют по две буквы.
\par Следовательно, любая гармоничная расстановка относится хотя бы к одному из двух случаев: либо каждая строка чередует две буквы, либо каждый столбец чередует две буквы. Действительно, если первый случай не выполнен, в некоторой строке встречаются хотя бы три буквы, и по доказанному выполнен второй случай.
\par Посчитаем первый случай. Множество букв первой строки выбирается $C_4^2=6$ способами. Затем множества букв остальных строк определены: соседние строки используют дополняющие друг друга пары. Для каждой из $2020$ строк независимо выбираем одну из двух букв, с которой начинается чередование. Получаем $6\cdot2^{2020}$ расстановок; все они гармоничны. По симметрии второй случай даёт столько же.
\par Дважды посчитаны расстановки, в которых и строки, и столбцы чередуют по две буквы. Они имеют период $2$ в обоих направлениях, поэтому полностью определяются левым верхним квадратом $2\times2$. Его можно заполнить четырьмя различными буквами $4!=24$ способами. Следовательно, ответ равен
\[
\boxed{12\cdot2^{2020}-24}.
\]\par

\item Имеются $12$ карандашей попарно различной длины. Сколькими способами можно уложить их в коробку в два слоя по шесть карандашей так, чтобы в каждом слое карандаши были упорядочены по возрастанию длины слева направо, а каждый карандаш верхнего слоя лежал строго над карандашом нижнего слоя и был короче него?
\par\smallskip\noindent\textbf{Решение.} Заменим длины карандашей их рангами $1,\ldots,12$. Верхний и нижний слои образуют таблицу $2\times6$ с возрастанием в строках и столбцах. По биекции со скобочными последовательностями из задачи 11а количество таких таблиц равно
\[
c_6=\frac{1}{7}C_{12}^6=132.
\]\par

\item Троих англичан, троих немцев и троих русских ставят в ряд. Сколько существует порядков, в которых никакие три соотечественника не стоят подряд? Все девять человек различны.
\par\smallskip\noindent\textbf{Решение.} Всего имеется $9!$ расстановок. Найдём число плохих расстановок, в которых хотя бы одна тройка соотечественников стоит подряд, по формуле включений-исключений.
\par Для одной фиксированной национальности начало блока из трёх человек можно выбрать $7$ способами, порядок внутри блока - $3!$ способами, а остальных людей расставить $6!$ способами. Сумма по трём национальностям равна $3\cdot7\cdot3!\cdot6!$.
\par Если подряд должны стоять тройки двух заданных национальностей, зафиксируем порядок этих двух блоков. Пусть $x_1,x_2,x_3$ - количества людей третьей национальности до первого блока, между блоками и после второго. Тогда $x_1+x_2+x_3=3$ в неотрицательных целых числах; решений $C_5^2=10$. Троих оставшихся людей можно упорядочить $3!$ способами, людей внутри каждого из двух блоков - ещё по $3!$ способами, а порядок самих блоков выбрать $2$ способами. Пар национальностей $3$, поэтому сумма попарных пересечений равна $3\cdot C_5^2\cdot2\cdot(3!)^3$.
\par Наконец, если все три национальности стоят отдельными блоками, имеется $(3!)^4$ расстановок: $3!$ порядков блоков и по $3!$ порядков людей внутри каждого. Следовательно, ответ равен
\[
\begin{aligned}
9!&-\left(3\cdot7\cdot3!\cdot6!
-3\cdot C_5^2\cdot2\cdot(3!)^3+(3!)^4\right)\\
&=362880-90720+12960-1296=283824.
\end{aligned}
\]\par

\item Сколько чисел от $1$ до $1\,000\,000$ включительно не являются ни квадратами, ни кубами, ни пятыми степенями целых чисел?
\par\smallskip\noindent\textbf{Решение.} Применим включения-исключения к множествам квадратов, кубов и пятых степеней. В указанном промежутке имеется $1000$ квадратов, $100$ кубов и $15$ пятых степеней. Их попарные пересечения состоят соответственно из шестых, десятых и пятнадцатых степеней; таких чисел $10$, $3$ и $2$. В тройном пересечении находятся тридцатые степени, и такая степень только одна - число $1$. Поэтому ответ равен
\[
1000000-1000-100-15+10+3+2-1=998899.
\]\par

\item На прямой отмечены $100$ точек. Сколькими способами можно выбрать из них $10$ точек так, чтобы никакие две выбранные точки не были соседними среди отмеченных?
\par\smallskip\noindent\textbf{Решение.} Пусть $x_1$ - число невыбранных точек до первой выбранной, $x_2,\ldots,x_{10}$ - числа невыбранных точек между соседними выбранными, а $x_{11}$ - число невыбранных точек после последней выбранной. Всего невыбранных точек $90$, поэтому
\[
x_1+x_2+\cdots+x_{11}=90,
\qquad x_1,x_{11}\ge0,\quad x_2,\ldots,x_{10}\ge1.
\]
Вычтем по единице из девяти внутренних промежутков. Получим уравнение $y_1+\cdots+y_{11}=81$ в неотрицательных целых числах. Его решения взаимно однозначно соответствуют допустимым выборам точек. По методу шаров и перегородок ответ равен
\[
C_{81+11-1}^{11-1}=C_{91}^{10}.
\]\par

\item Докажите, что число разбиений натурального числа $n$ на не более чем $m$ натуральных слагаемых равно числу его разбиений на натуральные слагаемые, каждое из которых не превосходит $m$. Порядок слагаемых не учитывается.
\par\smallskip\noindent\textbf{Решение.} Изобразим разбиение диаграммой из клеток: длины её строк равны слагаемым, записанным в невозрастающем порядке. Прочитаем ту же диаграмму по столбцам: высоты столбцов станут слагаемыми нового разбиения.
\par Если исходных слагаемых не более $m$, то каждый столбец имеет высоту не больше $m$, поэтому все новые слагаемые не превосходят $m$. Обратно, если все слагаемые не превосходят $m$, то столбцов не больше $m$, и после преобразования получится не более $m$ слагаемых. Повторное преобразование возвращает исходное разбиение; число клеток, то есть сумма слагаемых, сохраняется. Следовательно, построена биекция между требуемыми разбиениями.\par

\item Дан клетчатый квадрат $101\times101$. Внутри него выбирается квадрат $100\times100$, внутри выбранного - квадрат $99\times99$, и так далее до квадрата $1\times1$. Оказалось, что последний квадрат совпадает с центральной клеткой исходного квадрата. Сколько существует таких последовательностей квадратов? Ответ не должен содержать знака многоточия.
\par\smallskip\noindent\textbf{Решение.} Восстановим последовательность с конца: начнём с центральной клетки и за $100$ шагов расширим её до исходного квадрата. При каждом расширении нужно добавить одну строку сверху или снизу и один столбец слева или справа.
\par Обозначим числа расширений в этих направлениях через $U,D,L,R$. Изначально все счётчики равны нулю, а в конце каждый должен равняться $50$, поскольку центральная клетка отстоит от каждой стороны исходного квадрата на $50$ клеток. Каждый шаг однозначно задаётся парой: один выбор из $\{U,D\}$ и один из $\{L,R\}$.
\par Выберем $50$ из $100$ шагов, на которых увеличивается $U$; на остальных увеличивается $D$. Независимо выберем $50$ шагов, на которых увеличивается $R$; на остальных увеличивается $L$. Любые такие два выбора задают единственную последовательность квадратов. Все промежуточные квадраты остаются внутри исходного: ни один счётчик не превосходит своего конечного значения $50$. Ответ:
\[
\boxed{\left(C_{100}^{50}\right)^2}.
\]\par

\item На полоске из $n$ клеток фишка стоит в крайней левой клетке. За ход её можно передвинуть на одну или две клетки вправо. Докажите, что число способов добраться до крайней правой клетки равно $F_n$, где $F_1=F_2=1$ - первые числа Фибоначчи.
\par\smallskip\noindent\textbf{Решение.} Докажем утверждение по индукции. При $n=1$ имеется один способ - не делать ходов; при $n=2$ также один способ - сделать один ход. Это согласуется с $F_1=F_2=1$.
\par При $n\ge3$ последний ход имеет длину $1$ или $2$. В первом случае до него нужно добраться до предпоследней клетки, что по предположению индукции можно сделать $F_{n-1}$ способами; во втором - до клетки перед ней, что можно сделать $F_{n-2}$ способами. Случаи не пересекаются и исчерпывают все пути, поэтому их количество равно $F_{n-1}+F_{n-2}=F_n$.\par

\item Сколькими способами можно раскрасить клетки доски $20\times20$ в три цвета так, чтобы все цвета были использованы и в каждом уголке из трёх клеток присутствовали ровно два цвета?
\par\smallskip\noindent\textbf{Решение.} Сначала заметим, что в каждом квадрате $2\times2$ должны присутствовать ровно два цвета, каждый по два раза. Действительно, три одинаковые клетки дали бы одноцветный уголок, а при наличии трёх цветов нашёлся бы уголок с тремя разными цветами.
\par Рассмотрим сначала случай, когда две верхние левые клетки одинаковы. Тогда две клетки под ними также одинаковы между собой, но имеют другой цвет. Продолжая вниз, получаем, что первые два столбца совпадают, а их последовательные клетки имеют разные цвета. Цвет первой строки выбирается $3$ способами, каждой следующей - $2$ способами; всего $3\cdot2^{19}$ вариантов для первых двух столбцов.
\par Покажем, что для допустимой раскраски в этих столбцах обязательно используются все три цвета. Иначе цвета в них чередуются. В следующем столбце каждый вертикальный отрезок из двух клеток должен состоять из тех же двух цветов, поскольку вместе с двумя клетками предыдущего столбца он образует квадрат $2\times2$. Поэтому следующий столбец также чередует эти два цвета. Повторяя рассуждение, получаем, что на всей доске использованы только два цвета, что запрещено. Двухцветных чередующихся последовательностей ровно $3\cdot2=6$; их нужно исключить.
\par Если в первых двух столбцах используются все три цвета, то найдутся три последовательные строки с попарно различными цветами $a,b,c$. В самом деле, при отсутствии такой тройки каждый следующий цвет совпадал бы с позапрошлым, и вся последовательность была бы двухцветной. Рассмотрим клетки третьего столбца в этих трёх строках. Средняя из них должна иметь цвет из пары $\{a,b\}$ по верхнему квадрату $2\times2$ и из пары $\{b,c\}$ по нижнему. Значит, её цвет - $b$, а цвета соседних клеток - $a$ и $c$. Далее правило для квадратов $2\times2$ однозначно восстанавливает весь третий столбец: он совпадает с первыми двумя. Таким же образом восстанавливаются все остальные столбцы.
\par Итак, в этом случае все строки одноцветны, цвета соседних строк различны и используются все три цвета. Каждая из $3\cdot2^{19}-6$ последовательностей цветов строк действительно задаёт допустимую раскраску.
\par Симметрично, если верхняя левая клетка совпадает по цвету с клеткой под ней, получаем одноцветные столбцы и ещё $3\cdot2^{19}-6$ раскрасок. Эти два случая не пересекаются: иначе три клетки верхнего левого уголка были бы одного цвета.
\par Осталось проверить, что других случаев нет. Пусть верхняя левая клетка отличается и от правого, и от нижнего соседа. Тогда верхний левый квадрат имеет шахматный вид с двумя цветами. Его вертикальная сторона содержит два разных цвета; правило для квадратов $2\times2$ вынуждает каждый следующий столбец в первых двух строках содержать те же два цвета. Аналогично первые два столбца целиком используют только эти цвета. Теперь заполняем остальные клетки построчно: если три уже заполненные клетки квадрата $2\times2$ имеют только эти два цвета, четвёртая не может получить третий цвет, поскольку каждый цвет в квадрате должен встречаться дважды. Следовательно, вся доска была бы двухцветной. Этот случай недопустим.
\par Таким образом, ответ равен
\[
\boxed{2\bigl(3\cdot2^{19}-6\bigr)}.
\]\par

\item Фишка находится в левой нижней клетке доски $m\times n$. За ход её можно передвинуть на одну клетку вверх или вправо. Сколько существует путей в правую верхнюю клетку?
\par\smallskip\noindent\textbf{Решение.} Любой путь состоит из $m-1$ ходов вдоль одной стороны доски и $n-1$ ходов вдоль другой. Поэтому он однозначно задаётся последовательностью из $m+n-2$ ходов. Выберем $n-1$ позиций для ходов одного направления; остальные позиции займут ходы другого направления. Ответ:
\[
C_{m+n-2}^{n-1}.
\]\par

\item Сколько раз встречается каждая из цифр $0,1,\ldots,9$ при записи всех натуральных чисел от $1$ до $10^n-1$ без ведущих нулей?
\par\smallskip\noindent\textbf{Решение.} Зафиксируем ненулевую цифру $a$ и обозначим через $d(n)$ число её вхождений в записи натуральных чисел, содержащих не более $n$ разрядов. Положим $d(0)=0$.
\par Для подсчёта $d(n+1)$ разобьём записи на десять блоков по старшему разряду $b\in\{0,\ldots,9\}$. В каждом блоке младшие $n$ разрядов пробегают все записи от $0$ до $10^n-1$ и дают $d(n)$ вхождений цифры $a$. В блоке $b=a$ старший разряд даёт ещё $10^n$ вхождений. Временное добавление ведущих нулей не влияет на число вхождений ненулевой цифры, поэтому
\[
d(n+1)=10d(n)+10^n.
\]
Разделив на $10^{n+1}$, получаем
\[
\frac{d(n+1)}{10^{n+1}}=\frac{d(n)}{10^n}+\frac1{10}.
\]
С учётом $d(0)=0$ отсюда $d(n)/10^n=n/10$. Следовательно, каждая цифра $1,\ldots,9$ встречается $n10^{n-1}$ раз при $n\ge1$.
\par\textit{Цифра $0$.} Обозначим через $z_n$ число нулей в записях всех ровно $n$-значных чисел; $z_1=0$. Рассмотрим $(n+1)$-значные числа с фиксированной старшей цифрой $x\in\{1,\ldots,9\}$.
\par Если все остальные цифры нулевые, получаем одно число $x00\ldots0$, дающее $n$ нулей. Иначе пусть старший ненулевой разряд среди оставшихся $n$ разрядов имеет номер $k$ справа, где $1\le k\le n$. Последние $k$ цифр образуют произвольное $k$-значное число. Таких чисел $9\cdot10^{k-1}$; в их записях всего $z_k$ нулей. Между цифрой $x$ и этим числом стоят ещё $n-k$ нулей. Поэтому вклад данного случая при фиксированном $x$ равен
\[
9\cdot10^{k-1}(n-k)+z_k.
\]
Суммируя случаи и учитывая девять вариантов старшей цифры, получаем
\[
z_{n+1}=9\left[n+\sum_{k=1}^{n}\left(9\cdot10^{k-1}(n-k)+z_k\right)\right].
\]
Вычтем соответствующую формулу для $z_n$:
\[
z_{n+1}-z_n
=9\left[1+z_n+9\sum_{k=1}^{n-1}10^{k-1}\right]
=9z_n+9\cdot10^{n-1}.
\]
Следовательно, $z_{n+1}=10z_n+9\cdot10^{n-1}$. С начальным условием $z_1=0$ эта рекуррента даёт
\[
z_n=9(n-1)10^{n-2}\qquad(n\ge2).
\]
Суммируя по длинам записей, находим общее число нулей в требуемом промежутке:
\[
\sum_{k=1}^{n}z_k
=\sum_{k=2}^{n}9(k-1)10^{k-2}
=n10^{n-1}-\frac{10^n-1}{9}.
\]\par

\item Инверсией в перестановке называется пара чисел, в которой большее стоит раньше меньшего. Найдите суммарное число инверсий во всех перестановках чисел $1,2,\ldots,n$.
\par\smallskip\noindent\textbf{Решение.} При $n=1$ инверсий нет. Пусть $n\ge2$.
\par\textit{Первый способ.} Выберем два различных числа $a<b$. Они образуют инверсию ровно в половине всех $n!$ перестановок: обмен местами $a$ и $b$ задаёт биекцию между перестановками, в которых $a$ стоит раньше $b$, и перестановками с обратным порядком. Поэтому вклад каждой пары равен $n!/2$. Всего пар $C_n^2$, следовательно, суммарное число инверсий равно
\[
C_n^2\cdot\frac{n!}{2}=\frac{n(n-1)n!}{4}.
\]
\textit{Второй способ.} Сначала выберем пару чисел $a<b$, что можно сделать $C_n^2$ способами. Затем выберем две позиции для этих чисел - ещё $C_n^2$ способов. Чтобы пара образовывала инверсию, на более раннюю позицию нужно поставить $b$, а на более позднюю - $a$; этот порядок определяется однозначно. Оставшиеся $n-2$ числа расставляются $(n-2)!$ способами. Так каждая инверсия в каждой перестановке учитывается ровно один раз, поэтому ответ равен
\[
\left(C_n^2\right)^2(n-2)!=\frac{n(n-1)n!}{4}.
\]\par

\end{problems}
\end{document}
